% 约瑟夫·傅里叶（综述）
% license CCBYSA3
% type Wiki

本文根据 CC-BY-SA 协议转载翻译自维基百科\href{https://en.wikipedia.org/wiki/Joseph_Fourier}{相关文章}

\begin{figure}[ht]
\centering
\includegraphics[width=6cm]{./figures/78d7b0636c3dfd52.png}
\caption{} \label{fig_YSf_1}
\end{figure}
让-巴蒂斯特·约瑟夫·傅里叶（Jean-Baptiste Joseph Fourier，/ˈfʊrieɪ, -iər/；法语：[ʒɑ̃ batist ʒozɛf fuʁje]；1768年3月21日－1830年5月16日）是一位法国数学家和物理学家，出生于勃艮第的欧塞尔。他最为人熟知的是开创了傅里叶级数的研究，这一研究最终发展为傅里叶分析与调和分析，并广泛应用于热传导和振动问题的研究。傅里叶变换和傅里叶导热定律也因他而得名。傅里叶还通常被认为是温室效应的发现者。
\subsection{生平}
傅里叶出生于欧塞尔（今属法国约讷省），父亲是一名裁缝。他在九岁时成为孤儿。后来有人将他推荐给欧塞尔主教，经由这一引荐，他进入圣马克修道院的本笃会接受教育。军队科学部门的任职名额保留给出身高贵的人，而傅里叶因出身寒微不具备资格，因此他接受了一份数学军校讲师的职位。

在法国大革命期间，他在家乡地区积极参与革命活动，并在当地的革命委员会任职。他曾在恐怖统治时期被短暂监禁，但在1795年被任命为师范学校的教师，随后接替约瑟夫-路易·拉格朗日在综合理工学院的职位。

1798年，傅里叶作为拿破仑·波拿巴远征埃及时的科学顾问随行，并被任命为埃及学会的秘书。由于英国舰队的封锁，傅里叶组织了工坊，为法军提供所需的军火。他还为拿破仑在开罗创立的埃及学会（又称开罗学会）撰写了多篇数学论文，该学会旨在削弱英国在东方的影响。1801年，随着英国的胜利以及雅克-弗朗索瓦·梅努率领的法军投降，傅里叶返回法国。
\begin{figure}[ht]
\centering
\includegraphics[width=6cm]{./figures/9ee867e67b3cdc03.png}
\caption{} \label{fig_YSf_2}
\end{figure}
1801年，[4] 拿破仑任命傅里叶为伊泽尔省（格勒诺布尔）的省长（总督），由他负责监督道路建设及其他工程。然而，在此之前，傅里叶已从拿破仑的埃及远征中返回法国，准备恢复其在综合理工学院的教授职位，但拿破仑在一次评语中作出了不同的决定。

“……由于伊泽尔省的省长近日去世，我愿通过任命傅里叶公民担任此职来表达我对他的信任。[4]”
